% For compilation piecewise (see projekt.tex), it is necessary to uncomment it and change
% \documentclass[../projekt.tex]{subfiles}
% \begin{document}

%=========================================================================

\chapter{Introduction}
\label{introduction}

With the ever-increasing complexity of software systems and the growing demand for high-quality and secure code, formal verification methods have become essential in some parts of the software development industry \cite{qualitxinc}.
Formal verification methods are a set of mathematical techniques used to prove the correctness of software systems, going beyond traditional testing and debugging methods, which only test a finite number of inputs, allowing potentially critical errors to go undetected \cite{galois}.
The implementation of such formal methods is not straightforward, as it requires a deep understanding of both the verified software and the application of formal methods, typically necessitating specialized domain knowledge.
Formal verification tools, such as Predator, utilize static analysis to verify code written by developers.
In Predator's case, this involves verifying sequential C programs for memory safety while operating with pointers and linked lists \cite{predator}.
Predator is built on top of the Code Listener infrastructure \cite{code_listener}, which provides an application programming interface (API) foundation for creating such static analyzers.
However, the majority of the Code Listener infrastructure was written nearly fifteen years ago, which has led to certain challenges in continuing its development and maintenance in its current form.
To address these challenges, this bachelor's thesis focuses on refactoring the Code Listener infrastructure into a modern GCC \cite{gcc} plugin.

% TODO: motivation (why the refactor and the need for a new plugin, ...)

% TODO: solution approach (modularity, new C++ version, json, visualization, ...)

% TODO: end of introduction (structure of the thesis)

\chapter{Preliminaries}
\label{preliminaries}

preliminaries

\chapter{Design}
\label{design}

design

\chapter{Implementation}
\label{implementation}

implementation

\chapter{Testing}
\label{testing}

testing

\chapter{Conclusion}
\label{conclusion}

conclusion

%=========================================================================

% For compilation piecewise (see projekt.tex), it is necessary to uncomment it
% \end{document}